\documentclass[]{article}

\usepackage[margin=1in]{geometry}

\usepackage{amsmath,mathrsfs,amssymb,empheq,bm,pdfpages,cancel}

\usepackage[hidelinks]{hyperref}
\usepackage{titlesec}               %   paragraphs are like subsections
\setcounter{secnumdepth}{4}         %       Add section numbers
\setcounter{tocdepth}{4}            %       Add to the table of contents
\titleformat{\paragraph}[block]     %       Put a line break after the heading
  {\normalfont\normalsize\bfseries}{\theparagraph}{1em}{}
%   \subsection{\texorpdfstring{${\bm F}_{01}$}{F01}}

\newcommand{\pt}{\partial}

% \makeatletter                                   %  line spacing stretch for pmatrix,
% \renewcommand*\env@matrix[1][\arraystretch]{%   %     vmatrix, bmatrix.
%   \edef\arraystretch{#1}                        %  usage: \begin{pmatrix}[1.5]
%   \hskip -\arraycolsep                          %           A \\ B
%   \let\@ifnextchar\new@ifnextchar               %         \end{pmatrix}
%   \array{*\c@MaxMatrixCols c}}
% \makeatother
% 
% 
% \DeclareRobustCommand{\frac}[3][0pt]{%
%  {\begingroup\hspace{#1}#2\hspace{#1}\endgroup\over\hspace{#1}#3\hspace{#1}}}
%  %  fraction with a longer line.  usage: \frac[5pt]{}{}


\begin{document}

\title{XXX}
\maketitle

%=============================================================================
\section{Eq.(XXX)}  %<<<1

%>>>1

%===============    Folding   ===================
%<<<1
%<<<2
%>>>2
%>>>1

%===============    Include Figure   ===================

\begin{figure}[htb]
  \centering
  \includegraphics[width=18cm,trim={1cm 2cm 3cm 4cm},clip]{fig.pdf}
  \caption{
  CAPTION
  }
%
  \label{label}
\end{figure}


%===============    Font for code   ===================

\begin{verbatim}
>> Aval=np.array([12.04, 12.37, 12.35, 12.43, 12.34, 12.36, 12.48, 12.33, 12.33])
>> Bval=np.array([12.18, 12.37, 12.38, 12.36, 12.47, 12.48, 12.57, 12.28, 12.42])
>> diff = Aval-Bval
>> np.average(diff)
-0.05333333333333358
>> print(np.sum( (diff-np.average(diff))**2 ))
0.0498000000000005
\end{verbatim}

Or

\verb+Some letters+

%===============    Line Spacing Stretch for pmatrix, vmatrix, bmatrix.  ===================

\makeatletter                                   %  usage: \begin{pmatrix}[1.5]
\renewcommand*\env@matrix[1][\arraystretch]{%   %           A \\ B
  \edef\arraystretch{#1}                        %         \end{pmatrix}
  \hskip -\arraycolsep
  \let\@ifnextchar\new@ifnextchar
  \array{*\c@MaxMatrixCols c}}
\makeatother

%==============    Fraction with a longer line  =============================================

\DeclareRobustCommand{\frac}[3][0pt]{%
 {\begingroup\hspace{#1}#2\hspace{#1}\endgroup\over\hspace{#1}#3\hspace{#1}}}
 %  fraction with a longer line.  usage: \frac[5pt]{}{}


%==============    Equations with left curly bracket  =============================================

\begin{empheq}[left=\empheqlbrace]{align}
  a + b &= c, \\
  d + e &= f.
\end{empheq}

%==============    Matirx with indices =======================================
\usepackage{nicematrix}

\begin{align}
  A =
  \begin{pNiceMatrix}[first-row,first-col]
        &        &        &      j &        &      k &        &        \\
        &      0 & \cdots &      0 & \cdots &      0 & \cdots &      0 \\
        & \vdots &        & \vdots &        & \vdots &        & \vdots \\
     j  &      0 & \cdots &      0 & \cdots &      1 & \cdots &      0 \\
        & \vdots &        & \vdots &        & \vdots &        & \vdots \\
     k  &      0 & \cdots &      0 & \cdots &      0 & \cdots &      0 \\
        & \vdots &        & \vdots &        & \vdots &        & \vdots \\
        &      0 & \cdots &      0 & \cdots &      0 & \cdots &      0
  \end{pNiceMatrix}  
\end{align}

%=============================================================================

\end{document}

